% =====================================================================
%  chapters/minggu-07.tex
%  Bab 7 — CSS Layout: Grid & Responsive Design (Media Query)
% =====================================================================

\chapter{CSS Layout: Grid \& Responsive Design}
\label{chap:grid-responsive}

% ---------------------------------------------------------------------
% 1. CAPAIAN PEMBELAJARAN
% ---------------------------------------------------------------------
\begin{capaian}
Setelah menyelesaikan bab ini, mahasiswa diharapkan mampu:
\begin{itemize}
  \item menjelaskan konsep CSS \emph{grid} untuk \emph{layout} dua dimensi;
  \item membuat \emph{layout} \emph{grid} dengan baris dan kolom;
  \item menggunakan \emph{media query} untuk membuat halaman responsif;
  \item menerapkan prinsip \emph{mobile-first};
  \item membangun halaman yang tampil baik di layar besar maupun kecil.
\end{itemize}
\end{capaian}

% ---------------------------------------------------------------------
% 2. TUJUAN PRAKTIK
% ---------------------------------------------------------------------
\begin{tujuanpraktik}
Pada sesi praktik ini, mahasiswa akan:
\begin{itemize}
  \item menyiapkan folder proyek \texttt{latihan-grid};
  \item mempraktikkan properti \emph{grid} (\texttt{grid-template-columns}, \texttt{grid-template-rows}, \texttt{grid-column});
  \item membangun \emph{layout} halaman lengkap (header, sidebar, konten, footer);
  \item menerapkan \emph{media query} untuk tampilan responsif.
\end{itemize}
\end{tujuanpraktik}

% ---------------------------------------------------------------------
% 3. RINGKASAN TEORI
% ---------------------------------------------------------------------
\section{Pengantar CSS Grid}
\label{sec:pengantar-grid}

CSS Grid adalah sistem \emph{layout} \textbf{dua dimensi}---mengatur baris \textbf{dan} kolom sekaligus. Berbeda dengan \emph{flexbox} yang satu dimensi, \emph{grid} sangat cocok untuk \emph{layout} halaman penuh (header, sidebar, konten, footer).

\begin{lstlisting}[caption={Ilustrasi grid container dan cells}, label={lst:grid-illust}]
+--------------------------------------------------+
|  Grid Container (display: grid)                  |
|  +--------+  +--------+  +--------+             |
|  |  col 1 |  |  col 2 |  |  col 3 |   <- baris 1
|  +--------+  +--------+  +--------+             |
|  +--------+  +--------+  +--------+             |
|  |  col 1 |  |  col 2 |  |  col 3 |   <- baris 2
|  +--------+  +--------+  +--------+             |
+--------------------------------------------------+
\end{lstlisting}

\subsection{Media Query}
\label{subsec:media-query}

\emph{Media query} memungkinkan CSS \textbf{berubah} berdasarkan ukuran layar perangkat. Ini adalah kunci dari \emph{responsive design}.

\begin{lstlisting}[language=CSS, caption={Sintaks media query}, label={lst:media-query-basic}]
@media (max-width: 768px) {
    /* Aturan ini hanya berlaku saat layar <= 768px */
}
\end{lstlisting}

\subsection{Breakpoint Umum}
\label{subsec:breakpoint}

\begin{table}[h!]
\centering
\caption{Breakpoint umum untuk \emph{responsive design}}
\label{tab:breakpoint}
\begin{tabular}{|l|l|}
\hline
\textbf{Breakpoint} & \textbf{Target Perangkat} \\
\hline
\texttt{max-width: 600px} & Ponsel \\
\hline
\texttt{max-width: 768px} & Tablet \\
\hline
\texttt{max-width: 992px} & Laptop kecil \\
\hline
\texttt{min-width: 1200px} & Desktop besar \\
\hline
\end{tabular}
\end{table}

% ---------------------------------------------------------------------
% 4. PRAKTIKUM 1 — MENYIAPKAN PROYEK
% ---------------------------------------------------------------------
\section{Praktikum 1 --- Menyiapkan Proyek}
\label{sec:grid-praktik-1}

\begin{langkah}
  \item Buat folder \texttt{latihan-grid} di VSCode.
  \item Buat file \texttt{index.html}:
\end{langkah}

\begin{lstlisting}[language=HTML, caption={Struktur index.html}, label={lst:grid-index}]
<!DOCTYPE html>
<html lang="id">
<head>
    <meta charset="UTF-8">
    <meta name="viewport"
          content="width=device-width, initial-scale=1.0">
    <title>Latihan Grid</title>
    <link rel="stylesheet" href="style.css">
</head>
<body>
    <div class="grid">
        <div class="cell">1</div>
        <div class="cell">2</div>
        <div class="cell">3</div>
        <div class="cell">4</div>
        <div class="cell">5</div>
        <div class="cell">6</div>
    </div>
</body>
</html>
\end{lstlisting}

\begin{langkah}
  \item Buat file \texttt{style.css}:
\end{langkah}

\begin{lstlisting}[language=CSS, caption={Grid tiga kolom}, label={lst:grid-style}]
.grid {
    display: grid;
    grid-template-columns: 1fr 1fr 1fr;
    gap: 10px;
}

.cell {
    background-color: lightcoral;
    border: 2px solid darkred;
    padding: 30px;
    text-align: center;
    font-size: 20px;
}
\end{lstlisting}

\begin{langkah}
  \item Buka dengan Live Server. Enam sel tersusun dalam \textbf{3 kolom $\times$ 2 baris}.
\end{langkah}

% ---------------------------------------------------------------------
% 5. PRAKTIKUM 2 — MENGATUR KOLOM DAN BARIS
% ---------------------------------------------------------------------
\section{Praktikum 2 --- Mengatur Kolom dan Baris}
\label{sec:grid-praktik-2}

\begin{langkah}
  \item Ubah \texttt{grid-template-columns} menjadi berbagai variasi:
\end{langkah}

\begin{lstlisting}[language=CSS, caption={Kolom tidak sama lebar}, label={lst:grid-1fr2fr}]
.grid {
    display: grid;
    grid-template-columns: 1fr 2fr 1fr;
    gap: 10px;
}
\end{lstlisting}

\begin{langkah}
  \item Coba juga dengan satuan tetap:
\end{langkah}

\begin{lstlisting}[language=CSS, caption={Kolom tetap + fleksibel}, label={lst:grid-fixed}]
.grid {
    display: grid;
    grid-template-columns: 200px 1fr;
    gap: 10px;
}
\end{lstlisting}

\begin{langkah}
  \item Tambahkan \texttt{grid-template-rows}:
\end{langkah}

\begin{lstlisting}[language=CSS, caption={Mengatur tinggi baris}, label={lst:grid-rows}]
.grid {
    display: grid;
    grid-template-columns: 1fr 1fr 1fr;
    grid-template-rows: 100px 200px;
    gap: 10px;
}
\end{lstlisting}

\begin{langkah}
  \item Amati bagaimana tinggi baris berbeda.
\end{langkah}

\begin{tip}
\texttt{1fr} berarti ``1 fraksi dari ruang yang tersedia''. \texttt{1fr 2fr} berarti kolom kedua dua kali lebih lebar dari kolom pertama.
\end{tip}

\begin{tangkapanlayar}
  {assets/images/bab-07/grid-1fr2fr.png}
  {Hasil \texttt{index.html} dengan \texttt{grid-template-columns: 1fr 2fr 1fr} (kolom tengah lebih lebar).}
  {Ubah nilai \texttt{grid-template-columns}, simpan, dan amati perubahan di peramban.}
\end{tangkapanlayar}

% ---------------------------------------------------------------------
% 6. PRAKTIKUM 3 — MENEMPATKAN ITEM DI GRID
% ---------------------------------------------------------------------
\section{Praktikum 3 --- Menempatkan Item di Grid}
\label{sec:grid-praktik-3}

\begin{langkah}
  \item Buat file \texttt{layout.html}:
\end{langkah}

\begin{lstlisting}[language=HTML, caption={Struktur layout.html}, label={lst:layout-html}]
<!DOCTYPE html>
<html lang="id">
<head>
    <meta charset="UTF-8">
    <title>Layout Grid</title>
    <link rel="stylesheet" href="layout.css">
</head>
<body>
    <div class="layout">
        <header class="header">Header</header>
        <nav class="nav">Navigasi</nav>
        <main class="main">Konten Utama</main>
        <aside class="aside">Sidebar</aside>
        <footer class="footer">Footer</footer>
    </div>
</body>
</html>
\end{lstlisting}

\begin{langkah}
  \item Buat \texttt{layout.css}:
\end{langkah}

\begin{lstlisting}[language=CSS, caption={Layout halaman dengan grid}, label={lst:layout-css}]
.layout {
    display: grid;
    grid-template-columns: 1fr 3fr;
    grid-template-rows: auto auto 1fr auto;
    gap: 10px;
    min-height: 100vh;
}

.header {
    grid-column: 1 / -1;
    background-color: #2c3e50;
    color: white;
    padding: 20px;
}

.nav {
    background-color: #3498db;
    color: white;
    padding: 20px;
}

.main {
    background-color: #ecf0f1;
    padding: 20px;
}

.aside {
    background-color: #e67e22;
    color: white;
    padding: 20px;
}

.footer {
    grid-column: 1 / -1;
    background-color: #2c3e50;
    color: white;
    padding: 20px;
    text-align: center;
}
\end{lstlisting}

\begin{langkah}
  \item Buka di Live Server. Anda memiliki \emph{layout} halaman lengkap: header di atas, sidebar kiri, konten tengah, dan footer di bawah.
  \item Coba ubah \texttt{grid-template-columns} menjadi \texttt{1fr 2fr} dan amati perubahannya.
\end{langkah}

\begin{catatan}
\texttt{grid-column: 1 / -1} berarti item membentang dari kolom 1 sampai kolom terakhir.
\end{catatan}

\begin{tangkapanlayar}
  {assets/images/bab-07/layout.html.png}
  {Hasil \texttt{layout.html} (header, sidebar, konten, footer).}
  {Buka \texttt{layout.html} di peramban dan amati struktur \emph{layout} halaman.}
\end{tangkapanlayar}

% ---------------------------------------------------------------------
% 7. PRAKTIKUM 4 — PENGENALAN MEDIA QUERY
% ---------------------------------------------------------------------
\section{Praktikum 4 --- Pengenalan Media Query}
\label{sec:grid-praktik-4}

\begin{langkah}
  \item Buat file \texttt{responsive.html}:
\end{langkah}

\begin{lstlisting}[language=HTML, caption={Struktur responsive.html}, label={lst:responsive-html}]
<!DOCTYPE html>
<html lang="id">
<head>
    <meta charset="UTF-8">
    <meta name="viewport"
          content="width=device-width, initial-scale=1.0">
    <title>Responsive</title>
    <link rel="stylesheet" href="responsive.css">
</head>
<body>
    <h1>Halaman Responsif</h1>
    <div class="konten">
        <div class="kartu">Kartu 1</div>
        <div class="kartu">Kartu 2</div>
        <div class="kartu">Kartu 3</div>
    </div>
</body>
</html>
\end{lstlisting}

\begin{langkah}
  \item Buat \texttt{responsive.css}:
\end{langkah}

\begin{lstlisting}[language=CSS, caption={Grid responsif dengan media query}, label={lst:responsive-css}]
body {
    font-family: Arial, sans-serif;
    margin: 20px;
}

.konten {
    display: grid;
    grid-template-columns: 1fr 1fr 1fr;
    gap: 15px;
}

.kartu {
    background-color: #3498db;
    color: white;
    padding: 40px;
    text-align: center;
    border-radius: 8px;
}

/* Media query: saat layar <= 768px (tablet) */
@media (max-width: 768px) {
    .konten {
        grid-template-columns: 1fr 1fr;
    }
}

/* Media query: saat layar <= 480px (ponsel) */
@media (max-width: 480px) {
    .konten {
        grid-template-columns: 1fr;
    }
    h1 {
        font-size: 20px;
    }
}
\end{lstlisting}

\begin{langkah}
  \item Buka di Live Server. Persempit dan lebarkan jendela browser. Amati bagaimana jumlah kolom berubah: 3 $\rightarrow$ 2 $\rightarrow$ 1.
\end{langkah}

\begin{tip}
Gunakan DevTools (\texttt{F12}) $\rightarrow$ ikon perangkat (\emph{toggle device toolbar}) untuk mensimulasikan layar ponsel/tablet.
\end{tip}

\begin{tangkapanlayar}
  {assets/images/bab-07/responsive-desktop.png}
  {Hasil \texttt{responsive.html} di layar desktop (3 kolom).}
  {Buka \texttt{responsive.html} di peramban dengan lebar penuh.}
\end{tangkapanlayar}

\begin{tangkapanlayar}
  {assets/images/bab-07/responsive-mobile.png}
  {Hasil \texttt{responsive.html} di layar ponsel (1 kolom)---gunakan DevTools \emph{device toolbar}.}
  {Aktifkan DevTools $\rightarrow$ \emph{toggle device toolbar}, lalu pilih ukuran ponsel.}
\end{tangkapanlayar}

% ---------------------------------------------------------------------
% 8. PRAKTIKUM 5 — MOBILE-FIRST
% ---------------------------------------------------------------------
\section{Praktikum 5 --- Mobile-First}
\label{sec:grid-praktik-5}

Prinsip \textbf{mobile-first}: tulis CSS untuk layar kecil dulu, lalu gunakan \texttt{min-width} untuk layar yang lebih besar.

\begin{langkah}
  \item Buat file \texttt{mobile-first.css}:
\end{langkah}

\begin{lstlisting}[language=CSS, caption={Pendekatan mobile-first}, label={lst:mobile-first}]
/* Default: untuk ponsel (layar kecil) */
.konten {
    display: grid;
    grid-template-columns: 1fr;
    gap: 15px;
}

.kartu {
    background-color: #27ae60;
    color: white;
    padding: 40px;
    text-align: center;
    border-radius: 8px;
}

/* Layar >= 600px (tablet) */
@media (min-width: 600px) {
    .konten {
        grid-template-columns: 1fr 1fr;
    }
}

/* Layar >= 992px (desktop) */
@media (min-width: 992px) {
    .konten {
        grid-template-columns: 1fr 1fr 1fr;
    }
}
\end{lstlisting}

\begin{langkah}
  \item Buat file \texttt{mobile-first.html} yang memakai CSS ini (salin dari \texttt{responsive.html}, ganti nama file CSS).
  \item Buka di Live Server dan amati: di layar kecil tampil 1 kolom, makin lebar makin banyak kolom.
\end{langkah}

\begin{catatan}
\textbf{Perbedaan}: \emph{mobile-first} memakai \texttt{min-width} (dari kecil ke besar), sedangkan pendekatan biasa memakai \texttt{max-width} (dari besar ke kecil). Keduanya sah; \emph{mobile-first} lebih disarankan.
\end{catatan}

% ---------------------------------------------------------------------
% 9. PRAKTIKUM 6 — LAYOUT RESPONSIF LENGKAP
% ---------------------------------------------------------------------
\section{Praktikum 6 --- Membuat Layout Responsif Lengkap}
\label{sec:grid-praktik-6}

\begin{langkah}
  \item Buat file \texttt{halaman.html}:
\end{langkah}

\begin{lstlisting}[language=HTML, caption={Struktur halaman.html}, label={lst:halaman-html}]
<!DOCTYPE html>
<html lang="id">
<head>
    <meta charset="UTF-8">
    <meta name="viewport"
          content="width=device-width, initial-scale=1.0">
    <title>Halaman Responsif Lengkap</title>
    <link rel="stylesheet" href="halaman.css">
</head>
<body>
    <header class="header">
        <h1>Website Saya</h1>
        <nav>
            <ul class="menu">
                <li><a href="#">Beranda</a></li>
                <li><a href="#">Tentang</a></li>
                <li><a href="#">Kontak</a></li>
            </ul>
        </nav>
    </header>

    <main class="konten">
        <section class="artikel">
            <h2>Artikel Utama</h2>
            <p>Ini adalah konten utama halaman.
               Teks ini akan menyesuaikan lebar
               layar.</p>
        </section>
        <aside class="sidebar">
            <h3>Sidebar</h3>
            <p>Informasi tambahan.</p>
        </aside>
    </main>

    <footer class="footer">
        <p>&copy; 2026 Website Saya</p>
    </footer>
</body>
</html>
\end{lstlisting}

\begin{langkah}
  \item Buat \texttt{halaman.css} (mobile-first):
\end{langkah}

\begin{lstlisting}[language=CSS, caption={Gaya halaman responsif (mobile-first)}, label={lst:halaman-css}]
* {
    margin: 0;
    padding: 0;
    box-sizing: border-box;
}

body {
    font-family: Arial, sans-serif;
    line-height: 1.6;
}

.header {
    background-color: #2c3e50;
    color: white;
    padding: 20px;
    text-align: center;
}

.menu {
    list-style: none;
    display: flex;
    justify-content: center;
    gap: 20px;
    margin-top: 10px;
}

.menu a {
    color: white;
    text-decoration: none;
}

.konten {
    display: grid;
    grid-template-columns: 1fr;
    gap: 20px;
    padding: 20px;
    max-width: 1200px;
    margin: 0 auto;
}

.artikel {
    background-color: #ecf0f1;
    padding: 20px;
    border-radius: 8px;
}

.sidebar {
    background-color: #bdc3c7;
    padding: 20px;
    border-radius: 8px;
}

.footer {
    background-color: #2c3e50;
    color: white;
    text-align: center;
    padding: 15px;
}

/* Tablet: 2 kolom untuk konten */
@media (min-width: 600px) {
    .konten {
        grid-template-columns: 2fr 1fr;
    }
}

/* Desktop: menu di kanan header */
@media (min-width: 992px) {
    .header {
        display: flex;
        justify-content: space-between;
        align-items: center;
        text-align: left;
    }
    .menu {
        margin-top: 0;
    }
}
\end{lstlisting}

\begin{langkah}
  \item Buka di Live Server. Uji di berbagai ukuran layar (gunakan DevTools \emph{device toolbar}).
  \item Amati:
\end{langkah}

\begin{itemize}
  \item Di ponsel: konten 1 kolom, header di tengah.
  \item Di tablet: konten 2 kolom (artikel + sidebar).
  \item Di desktop: header jadi sejajar (logo kiri, menu kanan).
\end{itemize}

\begin{tangkapanlayar}
  {assets/images/bab-07/halaman-desktop.png}
  {Hasil \texttt{halaman.html} di layar desktop (konten 2 kolom, header sejajar).}
  {Buka \texttt{halaman.html} dengan lebar layar penuh.}
\end{tangkapanlayar}

\begin{tangkapanlayar}
  {assets/images/bab-07/halaman-mobile.png}
  {Hasil \texttt{halaman.html} di layar ponsel (konten 1 kolom, header di tengah).}
  {Aktifkan DevTools $\rightarrow$ \emph{device toolbar} dan pilih ukuran ponsel.}
\end{tangkapanlayar}

% ---------------------------------------------------------------------
% 10. LATIHAN MANDIRI
% ---------------------------------------------------------------------
\section{Latihan Mandiri}
\label{sec:grid-latihan}

\begin{latihan}[Portofolio Responsif]
Buat halaman portofolio yang responsif. Ketentuan:
\begin{itemize}
  \item Buat file \texttt{portofolio.html} dan \texttt{portofolio.css}.
  \item Halaman berisi header dengan nama Anda dan menu navigasi, bagian ``Tentang Saya'', bagian ``Proyek'' berisi minimal 6 kartu proyek, dan footer.
  \item Gunakan CSS \emph{grid} untuk menyusun kartu proyek.
  \item Gunakan \emph{media query} (\emph{mobile-first}) agar: ponsel 1 kolom, tablet 2 kolom, desktop 3 kolom.
  \item Gunakan \emph{flexbox} untuk menu navigasi.
\end{itemize}
\end{latihan}

Struktur HTML yang disarankan:

\begin{lstlisting}[language=HTML, caption={Kerangka portofolio.html}, label={lst:portofolio-kartu}]
<header class="header">
    <h1>Nama Anda</h1>
    <nav>
        <ul class="menu">
            <li><a href="#">Beranda</a></li>
            <li><a href="#">Tentang</a></li>
            <li><a href="#">Proyek</a></li>
        </ul>
    </nav>
</header>

<section class="tentang">
    <h2>Tentang Saya</h2>
    <p>...</p>
</section>

<section class="proyek">
    <h2>Proyek Saya</h2>
    <div class="grid-proyek">
        <div class="kartu">Proyek 1</div>
        <!-- ... 6 kartu ... -->
    </div>
</section>

<footer class="footer">
    <p>&copy; 2026 Nama Anda</p>
</footer>
\end{lstlisting}

% ---------------------------------------------------------------------
% 11. HASIL YANG DIHARAPKAN
% ---------------------------------------------------------------------
\section{Hasil yang Diharapkan}
\label{sec:grid-hasil}

Setelah menyelesaikan semua praktikum, Anda memiliki:
\begin{itemize}
  \item Folder \texttt{latihan-grid} berisi file: \path{index.html}, \path{style.css}, \path{layout.html}, \path{layout.css}, \path{responsive.html}, \path{responsive.css}, \path{mobile-first.html}, \path{mobile-first.css}, \path{halaman.html}, \path{halaman.css}.
  \item Halaman \texttt{index.html} dengan 6 sel dalam 3 kolom.
  \item Halaman \texttt{layout.html} dengan \emph{layout} lengkap: header, sidebar, konten, footer.
  \item Halaman \texttt{responsive.html} yang berubah dari 3 $\rightarrow$ 2 $\rightarrow$ 1 kolom saat layar dipersempit.
  \item Halaman \texttt{halaman.html} yang responsif dari ponsel hingga desktop.
  \item Halaman \texttt{portofolio.html} (latihan mandiri) dengan kartu proyek yang menyesuaikan jumlah kolom.
\end{itemize}

\begin{tangkapanlayar}
  {assets/images/bab-07/portofolio.html.png}
  {Hasil latihan mandiri \texttt{portofolio.html} di layar desktop (3 kolom kartu).}
  {Buat \texttt{portofolio.html} sesuai ketentuan latihan, lalu bandingkan tampilan Anda.}
\end{tangkapanlayar}

\begin{tangkapanlayar}
  {assets/images/bab-07/portofolio-mobile.png}
  {Hasil latihan mandiri \texttt{portofolio.html} di layar ponsel (1 kolom kartu).}
  {Gunakan DevTools $\rightarrow$ \emph{device toolbar} untuk menguji tampilan ponsel.}
\end{tangkapanlayar}

% ---------------------------------------------------------------------
% 12. TROUBLESHOOTING
% ---------------------------------------------------------------------
\section{Troubleshooting}
\label{sec:grid-trouble}

\begin{table}[h!]
\centering
\caption{Masalah umum CSS \emph{grid} dan \emph{responsive design}}
\label{tab:trouble-07}
\begin{tabular}{|p{3.5cm}|p{4cm}|p{5cm}|}
\hline
\textbf{Masalah} & \textbf{Penyebab} & \textbf{Solusi} \\
\hline
\emph{Grid} tidak tampil & \texttt{display: grid} belum ditambahkan & Pastikan properti \texttt{display: grid} ada di kontainer \\
\hline
Item tidak pindah kolom saat layar kecil & \emph{Media query} salah atau tidak ada & Pastikan ada \texttt{@media (max-width: ...)} atau \texttt{@media (min-width: ...)} \\
\hline
\emph{Media query} tidak berpengaruh & Urutan \emph{media query} salah & Letakkan \emph{media query} \textbf{setelah} aturan dasar; untuk \texttt{min-width}, urutkan dari kecil ke besar \\
\hline
\emph{Layout} berantakan di ponsel & Tidak ada \texttt{viewport} \emph{meta tag} & Pastikan ada \texttt{<meta name="viewport" content="width=device-width, initial-scale=1.0">} \\
\hline
Kolom tidak sama lebar & Menggunakan satuan tetap & Gunakan \texttt{1fr} untuk kolom yang fleksibel \\
\hline
Item menumpuk & \texttt{grid-template-columns} tidak sesuai & Periksa jumlah kolom yang didefinisikan \\
\hline
Perubahan tidak terlihat & Browser \emph{cache} & \emph{Refresh} dengan \texttt{Ctrl+Shift+R} \\
\hline
\end{tabular}
\end{table}

% ---------------------------------------------------------------------
% 13. RANGKUMAN
% ---------------------------------------------------------------------
\section{Rangkuman}
\label{sec:grid-rangkuman}

\begin{itemize}
  \item CSS \emph{grid} mengatur \emph{layout} \textbf{dua dimensi} (baris + kolom).
  \item \texttt{grid-template-columns} dan \texttt{grid-template-rows} menentukan struktur \emph{grid}.
  \item \texttt{1fr} = satu fraksi ruang yang tersedia.
  \item \texttt{grid-column: 1 / -1} membuat item membentang semua kolom.
  \item \emph{Media query} membuat halaman \textbf{responsif} terhadap ukuran layar.
  \item \textbf{Mobile-first}: tulis CSS ponsel dulu, lalu tingkatkan dengan \texttt{min-width}.
  \item Selalu sertakan \texttt{<meta name="viewport">} untuk halaman responsif.
\end{itemize}
